\section{Limitations}
\label{sec:limitations}
This study has several limitations which constrain its interpretability. First, the small sample size of five participants constrains the generalizability of the results. The standard deviations observed across the scores indicate high differences between individuals. 

Secondly, the spreadsheet knowledge level was relatively homogeneous. Only participants with spreadsheet knowledge between 1 and 6 took part in the study and the study did not assess how well expert-level users perform with the tool. 

Another important limitation is that the study did not have sessions in which the participants were only allowed to use traditional spreadsheet programs. This makes it impossible to know if the positive outcomes can be attributed to the flow graph visualization rather than other factors such as task familiarity, tasks that were too easy or simply the fact that the participants knew they were being observed.
Furthermore, the study did not impose any time constraints on the tasks nor was the correctness of the participants' answers evaluated. This constrains the interpretability of the NASA-TLX subscales: temporal demand and perceived performance.


Finally, the qualitative observations recorded during the study session could potentially introduce some degree of observational bias.